\section{Описание применения генеративной модели}

В рамках настоящей работы генеративные модели применялись в качестве вспомогательного инструмента на этапе проектирования пользовательского интерфейса дашборда. Основная цель их использования состояла в ускорении процесса создания и доработки макетов визуализаций, избавляя от необходимости ручного редактирования в графических редакторах.

Для задач редактирования существующих макетов и генерации новых прототипов на основе черновых рукописных эскизов использовалась мультимодальная генеративная модель \textbf{Nano Banana 2}, поддерживающая работу с изображениями в режиме диалога.

В процессе проектирования сформированные макеты вкладок дашборда загружались непосредственно в чат с моделью. Далее в текстовом сообщении формулировались конкретные правки, которые требовалось внести. Например:

\begin{itemize}
    \item Устранение визуальных аномалий от предыдущих итераций
    \item Изменение цветовой схемы
    \item Изменение типов и содержания графиков
\end{itemize}

Помимо редактирования готовых прототипов, модель Nano Banana 2 применялась для преобразования рукописных эскизов в оформленные макеты. Процесс состоял из следующих шагов: предварительно был нарисован от руки эскиз, отражающий желаемую компоновку элементов вкладки дашборда (расположение блоков, типы графиков, структура фильтров). Полученный эскиз был сфотографирован и загружен в диалог с моделью, после чего сформулирован промт: оформить эскиз в виде макета дашборда с соблюдением стандартов визуального представления данных.

Для подготовки макета одной из вкладок дашборда также использовалась генеративная модель \textbf{Recraft V4}, специализирующаяся на создании изображений по текстовому описанию.

В отличие от подхода с Nano Banana 2, где взаимодействие строилось на основе загружаемых изображений и итеративных правок, работа с Recraft V4 предполагала составление единого детального текстового промта. Промт содержал максимально подробное описание требуемого макета: перечень вкладок и их назначение, типы визуализаций для каждого блока (линейные графики, столбчатые диаграммы, индикаторы KPI, таблицы), расположение элементов на странице, цветовую схему, наличие и состав панели фильтров, а также стилистические требования к оформлению (деловой стиль, тёмная или светлая тема, типографика).

Оба инструмента дополняют друг друга: Nano Banana 2 обеспечивает гибкость редактирования на основе существующих изображений, тогда как Recraft V4 позволяет получить качественный первичный прототип по детальному текстовому описанию. Совместное применение этих моделей охватывает полный цикл дизайн-процесса — от первичного прототипирования до финальной доработки макетов вкладок дашборда.
